\begin{tikzpicture}[
    % Styles
    signal/.style={circle, draw=black!80, fill=blue!20, minimum size=6mm},
    conv_op/.style={rectangle, draw=black!80, fill=orange!30, minimum width=20mm, minimum height=10mm},
    residual/.style={diamond, draw=black!80, fill=green!30, minimum size=10mm},
    activation/.style={rectangle, draw=black!80, fill=yellow!30, minimum width=15mm, minimum height=8mm},
    % Arrows
    flow/.style={-latex, thick},
    skip/.style={-latex, thick, blue!60, dashed}
]

% Temporal Block Detail (showing one block)
\node[font=\large\bfseries] at (0, 6) {Temporal Block Internal Structure};

% Input
\node[signal] (input) at (0, 4) {$x$};
\node[below=2pt of input, font=\footnotesize] {[B, C, L]};

% First conv path
\node[conv_op] (conv1) at (0, 2) {Conv1d};
\node[activation] (act1) at (0, 0.5) {ReLU+Dropout};
\node[conv_op] (conv2) at (0, -1) {Conv1d};
\node[activation] (act2) at (0, -2.5) {ReLU+Dropout};

% Skip connection
\node[conv_op, minimum width=15mm] (skip) at (3, 1) {1×1 Conv};
\node[font=\footnotesize] at (3, 0.3) {(if needed)};

% Residual connection
\node[residual] (add) at (0, -4) {+};

% Output
\node[signal] (output) at (0, -5.5) {$y$};
\node[below=2pt of output, font=\footnotesize] {[B, C', L]};

% Connections
\draw[flow] (input) -- (conv1);
\draw[flow] (conv1) -- (act1);
\draw[flow] (act1) -- (conv2);
\draw[flow] (conv2) -- (act2);
\draw[flow] (act2) -- (add);

% Skip connection
\draw[skip] (input) -| (skip);
\draw[skip] (skip) |- (add);

\draw[flow] (add) -- (output);

% Chomp padding annotation
\draw[decorate, decoration={brace, amplitude=5pt, mirror}] 
    (4.5, 2.5) -- (4.5, -2.5) 
    node[midway, right=8pt, align=left, font=\footnotesize] {
        Chomp Padding\\removes future\\information
    };

% Full network overview on the right
\begin{scope}[shift={(10, 0)}]
    \node[font=\large\bfseries] at (0, 6) {Complete TCN Architecture};
    
    % Blocks
    \foreach \i/\d in {1/1, 2/2, 3/4} {
        \node[rectangle, draw=black!80, fill=orange!20, 
              minimum width=35mm, minimum height=20mm, rounded corners=3mm] 
              (block\i) at (0, 4-\i*2.2) {
            \begin{tabular}{c}
                \textbf{Temporal Block \i}\\
                \footnotesize dilation = \d\\
                \footnotesize 64 channels
            \end{tabular}
        };
    }
    
    % Input/Output
    \node[rectangle, draw=black!80, fill=green!30, 
          minimum width=35mm, minimum height=10mm, rounded corners=2mm] 
          (input_layer) at (0, 4.5) {Input (5 channels)};
    
    \node[rectangle, draw=black!80, fill=purple!30, 
          minimum width=35mm, minimum height=10mm, rounded corners=2mm] 
          (fc_layer) at (0, -3) {FC Layer (64→5)};
    
    \node[rectangle, draw=black!80, fill=red!30, 
          minimum width=35mm, minimum height=10mm, rounded corners=2mm] 
          (output_layer) at (0, -4.5) {Output (5 channels)};
    
    % Connections
    \draw[flow] (input_layer) -- (block1);
    \draw[flow] (block1) -- (block2);
    \draw[flow] (block2) -- (block3);
    \draw[flow] (block3) -- (fc_layer);
    \draw[flow] (fc_layer) -- (output_layer);
    
    % Annotations
    \node[align=center, font=\footnotesize] at (5, 0) {
        Total Receptive\\Field = $2^4 - 1 = 15$\\time steps
    };
\end{scope}

\end{tikzpicture}